\chapter{Foundations of Dynamical Systems}
\label{ch:foundations}

\epigraph{\itshape The whole of science is nothing more than a refinement of
everyday thinking.}{Albert Einstein}

\section{What is a dynamical system?}
A \emph{dynamical system} is a triple $(M,T,\Phi)$ where $M$ is a state
space (a manifold), $T$ is a time set (typically $\R$ or $\Z$), and
$\Phi:T\times M\to M$ is an evolution map satisfying the semi-group
axioms
\begin{align}
\Phi(0,\bm{x})&=\bm{x},\\
\Phi(t+s,\bm{x})&=\Phi(t,\Phi(s,\bm{x})).
\end{align}
Continuous-time systems are usually specified by an ODE
\begin{equation}
\dot{\bm{x}} = \bm{f}(\bm{x},t;\mu),\qquad \bm{x}(0)=\bm{x}_0,
\label{eq:ode}
\end{equation}
with parameter $\mu\in\R^p$. Discrete systems are given by a map
$\bm{x}_{k+1}=\bm{F}(\bm{x}_k)$. Spatially extended systems are modelled
by PDEs and become infinite-dimensional dynamical systems on
function spaces.

\begin{defn}{Flow and orbit}{flow}
The \emph{flow} of \cref{eq:ode} is the family of maps
$\{\Phi^t\}_{t\in\R}$ with $\Phi^t(\bm{x}_0)=\bm{x}(t)$. The
\emph{orbit} through $\bm{x}_0$ is the set
$\gamma(\bm{x}_0)=\{\Phi^t(\bm{x}_0):t\in\R\}$.
\end{defn}

\section{Linear systems and stability}
For the linear ODE $\dot{\bm{x}}=\bm{A}\bm{x}$ the solution is
$\bm{x}(t)=e^{t\bm{A}}\bm{x}_0$. The eigenvalues of $\bm{A}$ classify
the origin: stable node, unstable node, saddle, focus, centre. For
nonlinear systems, the \emph{Hartman--Grobman theorem} guarantees that,
near a hyperbolic fixed point, the flow is topologically conjugate to
its linearisation.

\begin{thm}{Lyapunov stability}{lyapunov}
Let $\bm{x}^\star$ be an equilibrium of \cref{eq:ode}. If there exists
a continuously differentiable function $V:\Omega\to\R$ on a neighbourhood
$\Omega\ni\bm{x}^\star$ with $V(\bm{x}^\star)=0$, $V(\bm{x})>0$ for
$\bm{x}\neq\bm{x}^\star$, and $\dot V(\bm{x})\le 0$ on $\Omega$, then
$\bm{x}^\star$ is Lyapunov stable. If $\dot V<0$ strictly, it is
asymptotically stable.
\end{thm}

\section{Nonlinear phenomena}
Nonlinearity introduces phenomena absent from linear theory:
\begin{itemize}
\item \textbf{Limit cycles:} isolated periodic orbits (Van der Pol).
\item \textbf{Bifurcations:} qualitative changes of phase portraits as
parameters vary --- saddle-node, transcritical, pitchfork, Hopf, period
doubling, homoclinic.
\item \textbf{Chaos:} sensitive dependence on initial conditions; the
maximal Lyapunov exponent
$\lambda_{\max}=\lim_{t\to\infty}\tfrac{1}{t}\log\|\delta\bm{x}(t)\|/\|\delta\bm{x}(0)\|$
is positive.
\item \textbf{Strange attractors:} fractal invariant sets attracting
nearby orbits (Lorenz, R\"ossler, Chua).
\end{itemize}

\begin{insight}{Why dynamics matters for ML}
Many learning algorithms are themselves dynamical systems: gradient
descent is an Euler discretisation of gradient flow $\dot\theta=-\nabla
L(\theta)$; momentum methods correspond to damped Hamiltonian
dynamics; sampling algorithms (Langevin, HMC) are stochastic flows on
parameter space. The mathematical machinery in this chapter is therefore
both the \emph{object} and the \emph{tool} of SciML.
\end{insight}

\section{Hamiltonian and Lagrangian mechanics}
Many physical systems possess a symplectic structure. Given a
configuration manifold $Q$ with coordinates $\bm{q}$ and conjugate
momenta $\bm{p}$, the \emph{Hamiltonian} $H(\bm{q},\bm{p},t)$ generates
flow via
\begin{equation}
\dot{\bm{q}}=\frac{\partial H}{\partial \bm{p}},\qquad
\dot{\bm{p}}=-\frac{\partial H}{\partial \bm{q}}.
\end{equation}
The 2-form $\omega=\dd\bm{q}\wedge\dd\bm{p}$ is preserved along the
flow. Equivalently, the Lagrangian $L(\bm{q},\dot{\bm{q}},t)=T-V$
satisfies the Euler--Lagrange equations
\begin{equation}
\frac{\dd}{\dd t}\frac{\partial L}{\partial \dot{\bm{q}}}-\frac{\partial L}{\partial \bm{q}}=0.
\end{equation}
Noether's theorem connects continuous symmetries to conserved
quantities --- a principle that motivates both equivariant neural
architectures (\cref{ch:equivariant}) and Hamiltonian/Lagrangian neural
networks (\cref{ch:nodes}).

\section{Partial differential equations}
A PDE is a relation
$\Nop[u](\bm{x},t)=0$ on a domain $\Omega\subset\R^d$, with boundary
conditions on $\partial\Omega$ and (for evolution problems) initial
conditions at $t=0$. Major classes:
\begin{itemize}
\item \textbf{Elliptic:} $-\Delta u=f$ (Poisson), steady-state.
\item \textbf{Parabolic:} $\partial_t u=\Delta u+f$ (heat), diffusive.
\item \textbf{Hyperbolic:} $\partial_t^2 u=c^2\Delta u$ (wave),
finite propagation speed.
\item \textbf{Conservation laws:} $\partial_t u+\nabla\cdot\bm{F}(u)=0$
(Burgers, Euler, Navier--Stokes).
\end{itemize}

Well-posedness in the sense of Hadamard requires existence,
uniqueness, and continuous dependence on data. For nonlinear hyperbolic
laws, \emph{weak} (entropy) solutions handle shock formation; we will
return to this when discussing PINN failure modes
(\cref{sec:pinn-failure}).

\section{Function spaces and the operator picture}
Solutions of PDEs live in function spaces (Sobolev, Bochner, etc.).
The map from problem data (coefficients, forcing, BCs, ICs) to solution
\begin{equation}
\Gop:\X\to\Y,\qquad a\mapsto u
\end{equation}
is the \emph{solution operator}. Most modern SciML --- especially
neural operators (\cref{ch:operators}) --- aims to learn $\Gop$
directly, rather than its values $u$ for fixed $a$.

\begin{recipe}{Picking the right viewpoint}
\begin{enumerate}
\item Single trajectory, fixed parameters: ODE/PDE \emph{solver}.
\item Family of trajectories, varying ICs: \emph{flow map} learning
(neural ODE, time-stepping operator).
\item Family of problem instances: \emph{solution operator} learning
(DeepONet, FNO).
\item Sparse data, unknown equation: \emph{system identification}
(SINDy, UDEs).
\end{enumerate}
\end{recipe}

\section{Stochastic dynamics}
When noise matters, replace the ODE by an SDE
\begin{equation}
\dd\bm{X}_t = \bm{f}(\bm{X}_t,t)\dd t + \bm{\sigma}(\bm{X}_t,t)\dd \bm{W}_t,
\end{equation}
with $\bm{W}_t$ a Wiener process. The Fokker--Planck equation governs
the density:
\begin{equation}
\partial_t p = -\nabla\!\cdot(\bm{f}p) + \tfrac{1}{2}\nabla\nabla:(\bm{\sigma}\bm{\sigma}^\T p).
\end{equation}
This connection between trajectories and densities is central to
\emph{score-based diffusion models} (\cref{ch:generative}), where the
reverse-time SDE is parameterised by a learned score
$\nabla\log p_t$.

\section{Inverse problems}
Many SciML applications are inverse: given observations
$\bm{y}=\Hop\bm{u}+\bm{\eta}$ of a state governed by a PDE, infer
parameters or initial conditions. Tikhonov regularisation, Bayesian
inference, and adjoint methods provide the classical scaffolding;
deep learning provides priors (deep image prior, score priors) and
amortised solvers.

\section{Summary}
Dynamical systems theory provides the language of states, flows,
stability, and conservation. Numerical analysis (\cref{ch:numerics})
provides tools to compute. Modern SciML threads neural networks
through both: as solvers, as surrogates, as identifiers, and as priors.
